\section{Exact Fej\'er resolution and hard-band asymptotics}
\label{sec:band-corollaries}

The sharp interval in \cref{thm:fixed-ray-mean-square} leaves an
$O(X\cE)$ boundary error.  A kernel whose Fourier transform has compact
support removes that error exactly.

Define the normalized Fej\'er kernel
\begin{equation}
 \kappa(t):=\frac2\pi\left(\frac{\sin(t/2)}t\right)^2,
 \qquad \kappa(0):=\frac1{2\pi}.
 \label{eq:normalized-Fejer-kernel}
\end{equation}
With the convention
$\widehat\kappa(\xi)=\int_\R\kappa(t)e^{-it\xi}dt$, one has
\begin{equation}
 \widehat\kappa(\xi)=(1-|\xi|)_+,
 \qquad \int_\R\kappa(t)dt=1.
 \label{eq:Fejer-transform}
\end{equation}

\begin{theorem}[Exact factor-ray Parseval identity]
\label{thm:exact-Fejer-ray-Parseval}
If $T\delta_X\geq1$, where $\delta_X$ is defined in
\eqref{eq:exact-ratio-spacing}, then
\begin{equation}
 \int_\R\kappa(t/T)|\cR_{h,D}(X;t)|^2dt
 =T\cE_{h,D}(X).
 \label{eq:exact-Fejer-ray-Parseval}
\end{equation}
The same identity holds for the raw transform and raw ray energy.
\end{theorem}

\begin{proof}
Insert \eqref{eq:ray-compression}, expand the square, and use
\eqref{eq:Fejer-transform}:
\[
 \int_\R\kappa(t/T)e^{it(\lambda_\rho-\lambda_{\rho'})}dt
 =T\widehat\kappa\!\left(T(\lambda_{\rho'}-\lambda_\rho)\right).
\]
For distinct rays the argument has absolute value at least
$T\delta_X\geq1$, so the triangle function vanishes, including at the
endpoints.  Every diagonal term has multiplier $T$.
\end{proof}

Thus an effective bandwidth $T\asymp_W X$ exactly resolves all primitive
factor rays.  It does not resolve different $k$ inside one ray; those terms
have already been combined into $c_{a,b}$.

\begin{corollary}[Averaged-shift Fej\'er-band asymptotic]
\label{cor:averaged-Fejer-band}
Under the hypotheses of \cref{thm:averaged-ray-energy}, suppose also that
$T\delta_X\geq1$.  Then
\begin{equation}
 \begin{aligned}
 &\sum_hG(h/H)
 \int_\R\kappa(t/T)|\cT_{h,D}(X;t)-\cM_D(X;t)|^2dt\\
 &\quad=\frac12THX I_0(G)I_2(W)\log H
 \bigl((\log X)^2-(\log D)^2\bigr)\\
 &\qquad\quad+O_{\eps,C,W,G}\!\left(
 THX(\log(2X))^2\log\log(3X)\right).
 \end{aligned}
 \label{eq:averaged-Fejer-band}
\end{equation}
If $G\not\equiv0$ and $W\not\equiv0$, this is an asymptotic throughout
the stated $D,H$ range.
\end{corollary}

\begin{proof}
Sum \eqref{eq:exact-Fejer-ray-Parseval} with weight $G(h/H)$ and insert
\eqref{eq:averaged-ray-energy}.  Since
$D\leq X^{1-\eps}$ and $H\geq X$, the main term is
$\gg_{\eps,G,W}THX(\log X)^3$ when $I_0(G)I_2(W)>0$, whereas the error is
smaller by a factor $O(\log\log X/\log X)$.
\end{proof}

\begin{corollary}[Averaged-shift hard-band formula]
\label{cor:averaged-hard-band}
Under the hypotheses of \cref{thm:averaged-ray-energy}, for every
$T\geq1$,
\begin{equation}
 \begin{aligned}
 &\sum_hG(h/H)
 \int_{-T}^{T}|\cT_{h,D}(X;t)-\cM_D(X;t)|^2dt\\
 &\quad=THX I_0(G)I_2(W)\log H
 \bigl((\log X)^2-(\log D)^2\bigr)\\
 &\qquad\quad+O_{\eps,C,W,G}\!\left(
 THX(\log(2X))^2\log\log(3X)
 +HX^2(\log(2X))^3\right).
 \end{aligned}
 \label{eq:averaged-hard-band}
\end{equation}
Consequently \eqref{eq:averaged-hard-band} is an asymptotic whenever
$T/X\to\infty$.
\end{corollary}

\begin{proof}
Sum \eqref{eq:hard-band-ray-law} over $h$.  Twice the leading term in
\eqref{eq:averaged-ray-energy} gives the displayed main term.  The first
error is twice $T$ times the averaged ray-energy error.  For the
Montgomery--Vaughan boundary error, \cref{thm:averaged-ray-energy} and its
proof give
\[
 \sum_hG(h/H)\cE_{h,D}(X)\ll_{\eps,C,W,G}HX(\log X)^3.
\]
Multiplication by $X$ gives the second error in
\eqref{eq:averaged-hard-band}.  Relative to the main term, the two errors
are respectively $O(\log\log X/\log X)$ and $O(X/T)$.
\end{proof}

\begin{corollary}[The TPC--10 truncation]
\label{cor:TPC10-truncation}
Let
\[
 D=X^{1/2}(\log X)^{-B}
\]
with fixed $B>0$.  Then
\begin{equation}
 (\log X)^2-(\log D)^2
 =\left(\frac34+o(1)\right)(\log X)^2.
 \label{eq:three-quarter-source-energy}
\end{equation}
Thus the leading coefficient contributed by the source cutoff is
$3/8$ in \eqref{eq:averaged-Fejer-band} and $3/4$ in
\eqref{eq:averaged-hard-band}, before the common factors
$THX I_0(G)I_2(W)\log H(\log X)^2$ are restored.
\end{corollary}

\begin{remark}[Raw target]
By the raw part of \cref{thm:averaged-ray-energy}, both band corollaries
remain valid with $\cT_{h,D}$ in place of $\cT_{h,D}-\cM_D$.  Centering is
used because $\cM_D$ is the long-shift first-moment model, not because it
changes the leading second-moment scale.
\end{remark}

\subsection{A parity-conditioned even-shift ensemble}

The all-shift average above is not forced by the proof.  To isolate the
local parity factor without selecting an individual shift, define
\begin{equation}
 \cM_D^{(2)}(X;t)
 :=2\sum_{\substack{n>D,\ m\geq1\\2\nmid mn}}
 u(n)W(mn/X)(n/m)^{it}
 \label{eq:even-shift-model}
\end{equation}
For even $h$, put $\cR^{(2)}_{h,D}=\cT_{h,D}-\cM_D^{(2)}$ and define
\begin{equation}
 \eta_r^{(2)}(h):=\Lambda(r+h)-2\one_{2\nmid r}.
 \label{eq:even-target-residual}
\end{equation}
Then the ray coefficient of $\cR^{(2)}_{h,D}$ is
\begin{equation}
 c^{(2)}_{a,b}(h;X,D)
 =\sum_{\substack{k\geq1\\ak>D}}
 u(ak)\eta_{abk^2}^{(2)}(h)W(abk^2/X),
 \label{eq:even-ray-coefficient}
\end{equation}
and $\cE^{(2)}_{h,D}$ denotes the sum of its squared moduli.  The factor
$2$ in \eqref{eq:even-shift-model} is fixed by the first-moment identity
\begin{equation}
 \frac2H\sum_{h\equiv0\ (2)}G(h/H)\cT_{h,D}(X;t)
 =I_0(G)\cM_D^{(2)}(X;t)
 +O_{A,C,W,G}\!\left(X(\log X)^{-A}\right),
 \label{eq:even-first-moment}
\end{equation}
uniformly for real $t$ and $X\leq H\leq X^C$.

\begin{proposition}[Even-shift ensemble]
\label{prop:even-shift-ensemble}
Under the hypotheses of \cref{thm:averaged-ray-energy}, the ray energy
$\cE^{(2)}_{h,D}$ of $\cR^{(2)}_{h,D}$ satisfies
\begin{equation}
 \begin{aligned}
 \sum_{h\equiv0\ (2)}G(h/H)\cE^{(2)}_{h,D}(X)
 ={}&\frac14HX I_0(G)I_2(W)\log H\\
 &\times\bigl((\log X)^2-(\log D)^2\bigr)\\
 &+O_{\eps,C,W,G}\!\left(
 HX(\log X)^2\log\log(3X)\right).
 \end{aligned}
 \label{eq:even-shift-ray-energy}
\end{equation}
Consequently the hard-band leading coefficient is one half of that in
\eqref{eq:averaged-hard-band}, and the Fej\'er-band leading coefficient is
one half of that in \eqref{eq:averaged-Fejer-band}.
\end{proposition}

\begin{proof}
For $r\in[\alpha X,\beta X]$, the prime number theorem modulo $2$ gives
\begin{equation}
 \sum_{h\equiv0\ (2)}G(h/H)\Lambda(r+h)
 =H I_0(G)\one_{2\nmid r}+O_A(H(\log H)^{-A}).
 \label{eq:even-target-first-moment}
\end{equation}
Summing this estimate against the source rows proves
\eqref{eq:even-first-moment}.  The same argument for $\Lambda^2$, followed
by expansion of \eqref{eq:even-target-residual}, gives
\begin{equation}
 \sum_{h\equiv0\ (2)}G(h/H)|\eta_r^{(2)}(h)|^2
 =H I_0(G)\log H\,\one_{2\nmid r}+O_{C,W,G}(H).
 \label{eq:even-target-second-moment}
\end{equation}
Indeed, for odd $r$ the conditional density of primes in the odd residue
class compensates for sampling half of the shifts.  For even $r$, the
target is even and $\Lambda(r+h)$ is supported only on powers of two, so
its contribution is absorbed by the error.

For distinct $r,r'\in[\alpha X,\beta X]$, one also has
\begin{equation}
 \sum_{h\equiv0\ (2)}G(h/H)
 |\eta_r^{(2)}(h)\eta_{r'}^{(2)}(h)|
 \ll_{C,W,G}H\log\log(3X).
 \label{eq:even-two-target-upper-sieve}
\end{equation}
If $r$ and $r'$ are odd, write $h=2j$ and apply the two-linear-forms
Selberg upper sieve to $2j+r$ and $2j+r'$.  If either is even, its target
factor is supported on powers of two; the resulting polylogarithmic
contribution is smaller than the displayed bound.  This also covers the
mixed-parity rows, which are present in \eqref{eq:even-ray-coefficient}.

The odd-product source diagonal is
\begin{equation}
 \sum_{\substack{n>D,\ m\geq1\\2\nmid mn}}
 u(n)^2|W(mn/X)|^2
 =\frac14X I_2(W)
 \bigl((\log X)^2-(\log D)^2\bigr)+O_W(X\log X).
 \label{eq:even-source-diagonal}
\end{equation}
Indeed, the odd $m$ sum contributes one half of the continuous density,
whereas removing the prime $2$ does not change the leading
$\frac12\log^2 y$ term in the weighted $u(n)^2$ sum over odd $n$.

Expand the squared ray coefficients.  Equations
\eqref{eq:even-target-second-moment} and
\eqref{eq:even-source-diagonal} give the displayed main term; even target
diagonal rows contribute only $O(H\cS_D(X))$.  For $k\ne\ell$, apply
\eqref{eq:even-two-target-upper-sieve} and then the full source ledger
\eqref{eq:same-ray-source-ledger}.  The resulting error is
$O(HX(\log X)^2\log\log(3X))$, proving
\eqref{eq:even-shift-ray-energy}.  The band statements follow from the
same factor-ray identities.
\end{proof}

\begin{remark}[No individual-shift selection]
The ensemble in \cref{prop:even-shift-ensemble} contains only admissible
parity shifts, but its exceptional behavior may still be concentrated at
any prescribed even shift.  No individual shift is selected.
\end{remark}
